\chapter{Conclusioni}
\label{chap:conclusioni}

\section{Risultati raggiunti}

Il \autoref{chap:introduzione} affida a RAM-USB l'obiettivo di progettare, implementare e validare un
servizio di backup remoto multi-utente rispetto ai requisiti della SRS \cite{ramusb-srs}, affidandosi ai
principi zero-trust e defense-in-depth per difendere il sistema da attacchi esterni, non per competere
con le soluzioni commerciali esistenti ma per discuterne le scelte di design in modo trasparente.

I dieci requisiti utente elencati nel \autoref{chap:requisiti} sono, con una sola eccezione, coperti da requisiti di
sistema implementati e tracciati fino al codice. Un utente può registrarsi fornendo solo email, password
e chiave pubblica SSH (RU-01), autenticarsi nelle sessioni successive (RU-02), caricare e recuperare
backup cifrati lato client da qualunque luogo (RU-03, RU-07), farlo solo dopo essersi autenticato (RU-04),
con la garanzia che né la password né il contenuto dei backup siano mai leggibili da un amministratore
(RU-05, RU-06), in uno spazio isolato da quello degli altri utenti (RU-08). Un amministratore può
osservare la salute e le prestazioni del sistema senza mai toccare i dati dei singoli utenti (RU-10).

L'unica eccezione è RU-09, la garanzia che nessuno possa modificare o cancellare un backup: è stata
dichiarata fuori perimetro fin dalla SRS, per i tempi stretti del progetto piuttosto che per una scelta di
design, ed è ripresa nella prossima sezione.

Le cinque sequenze d'uso che attraversano il sistema, registrazione, login, backup, restore e
consultazione delle metriche, non sono rimaste un esercizio sulla carta: sono state eseguite end-to-end
sul deployment reale su Proxmox, non solo sullo stack Docker di sviluppo, con gli stessi risultati
riportati nel \autoref{chap:testing}.

\section{Limiti noti}

La SRS accetta esplicitamente quattro rischi come costo del rilascio nei tempi previsti, non come
sviste \cite{ramusb-srs}. Dichiararli qui, uno per uno, è la stessa disciplina di onestà già applicata nel
\autoref{chap:sicurezza}, dove ogni proprietà che regge viene dichiarata insieme a quella che cade.

RISK-01 è RU-09 stesso: nessun requisito di sistema impedisce a un amministratore di modificare o
cancellare un backup. Resta un miglioramento desiderabile, non ancora costruito.

RISK-02 riguarda la master key: risiede in una variabile d'ambiente, senza una procedura di backup
(DV-F-18) né di rotazione (DV-F-19), entrambe requisiti "dovrebbe" e non "deve". La sua perdita
comporterebbe la perdita irreversibile dell'accesso a tutti i dati cifrati; la sua compromissione
romperebbe la garanzia zero-knowledge per ogni utente.

RISK-03 riguarda il client, pensato per girare nativamente sulla macchina dell'utente e non come
container Docker. Containerizzarlo è stato considerato e scartato per ora: un container non vede percorsi
arbitrari dell'host, come il Desktop dell'utente, a meno di un bind mount esplicito, e l'insieme dei file
da salvare si sceglie liberamente al momento del backup, non è noto in anticipo come per ogni altro
componente. L'opzione che perde meno isolamento, se il tema fosse ripreso, è un bind mount per singola
invocazione, limitato alla sola cartella oggetto del backup.

RISK-04 è la dipendenza circolare della risoluzione DNS di Network-Manager dal proprio stesso MagicDNS: se
Network-Manager pubblicasse una policy ACL non funzionante, potrebbe perdere anche la capacità di
risolvere l'hostname di Headscale di cui avrebbe bisogno per correggerla. La scelta di accettare comunque
il MagicDNS, invece di costruire un percorso di recupero automatico, è un rischio accettato per scelta
esplicita, non un difetto scoperto in seguito: in questo scenario di guasto, un operatore interviene
manualmente.

\section{Sviluppi futuri}

\subsection{Mitigazione dell'enumerazione degli indirizzi email}

Il \autoref{sec:scenari-attacco} analizza RISK-05: la registrazione rivela, tramite il codice HTTP
restituito, se un indirizzo è già registrato. Il comportamento ideale risponderebbe in modo identico in
entrambi i casi, per esempio sempre con un 202, mentre l'esito reale viaggerebbe fuori banda, in un'email
di conferma inviata all'indirizzo indicato: solo chi possiede davvero quella casella riceverebbe il
seguito, mentre chi sonda l'endpoint riceverebbe sempre la stessa risposta. Un effetto collaterale
positivo scomparirebbe insieme al problema: senza un account attivato dalla conferma, sparirebbe anche
l'occupazione abusiva di un indirizzo altrui.

Il costo non va taciuto: servirebbe una capacità di invio email che il sistema oggi non ha,
un client SMTP, token di conferma con scadenza, e una deliverability affidabile. Cambierebbe inoltre
UC-01, perché la pre-auth key oggi viaggia dentro la risposta 201 della registrazione, mentre con una
registrazione in due tempi dovrebbe essere consegnata solo dopo la conferma.

\subsection{Le due regole della policy delle password}

Il \autoref{sec:entropia} misura che, sotto scelta uniforme, il vincolo di almeno tre categorie di
carattere vale meno di un decimo di bit di entropia, mentre un carattere aggiuntivo ne vale 6,57: sotto
quell'ipotesi la regola di composizione è ridondante per costruzione, perché una password davvero casuale
usa già più categorie da sé. Il suo scopo reale, però, non è aumentare l'entropia, ma vincolare la scelta
di una persona, e lì il suo valore non è misurabile senza la distribuzione reale delle password scelte
dagli utenti.

Restano due direzioni da confrontare: tenere entrambe le regole così come sono, oppure eliminare il
vincolo di composizione e alzare la lunghezza minima, la strada indicata dalle linee guida NIST correnti
\cite{nist-sp-800-63b}, che raccomandano ai verificatori di non imporre regole di composizione e portano
il minimo per una password usata da sola a 15 caratteri. Anche la seconda strada ha però un costo da non
nascondere: senza una lista di password compromesse note, togliere il vincolo di composizione rimuove
una protezione debole senza sostituirla con una forte, e una password come \texttt{password} tornerebbe
accettabile.

\subsection{Hash con chiave per la ricerca dell'email}

Il \autoref{sec:scenari-attacco} analizza anche RISK-06: la chiave di ricerca dell'email è oggi uno
SHA-256 senza chiave, veloce da calcolare su uno spazio di indirizzi piccolo e indovinabile. Il
comportamento ideale deriverebbe quella chiave con una funzione dotata di chiave, HMAC-SHA256 sotto il
pepper, e non con la concatenazione ingenua che nel percorso della password basta solo perché lì
interviene Argon2id. Un pepper è compatibile con la ricerca, a differenza di un salt: essendo un unico
segreto globale, il digest resta deterministico e la ricerca al login continua a funzionare identica. Il
guadagno sarebbe che una copia del database, da sola, non permetterebbe più di verificare se un indirizzo
sia registrato, la stessa protezione che Argon2id offre già per la password. Non
risolverebbe invece RISK-05, che resta un attacco online contro l'API: sono due superfici diverse.

Il costo dichiarato dalla SRS per questa modifica presuppone però una reversibilità che oggi non esiste
ancora nel codice, e va corretto qui. La SRS descrive ogni \texttt{email\_hash} come ricostruibile
decifrando l'email con la master key, dato che è cifrata e non solo hashata; il record salvato, insieme
alla master key, contiene sì in linea di principio tutto il necessario.

\texttt{encryption.go}, però, dichiara oggi solo \texttt{EncryptEmail}, non la funzione inversa: nessun
percorso del codice decifra mai quella colonna. La reversibilità è quindi, per ora, una proprietà dei dati
salvati, non ancora del sistema: la migrazione dovrebbe prima scrivere quel percorso di decifratura
mancante, non solo invalidare e ricalcolare gli hash esistenti.

Il costo resta comunque reale anche a percorso scritto: il pepper diventerebbe portante anche per la
ricerca dei record, non solo per la verifica delle password, spostando la conseguenza della sua perdita da
"non si verificano più le password" a "non si trovano più i record", da valutare accanto a RISK-02.

\subsection{Altri requisiti "dovrebbe" non implementati}

Oltre a DV-F-18 e DV-F-19, già discussi come RISK-02, restano "dovrebbe" e non ancora costruiti il limite
di quota per utente (ST-F-14) e la replica automatica dei backup su CephFS (ST-F-15), entrambi rimandati
nella SRS come miglioramenti desiderabili piuttosto che vincoli del progetto. Resta inoltre da fare, sul
piano della verifica piuttosto che dell'implementazione, un test di carico e di stress che misuri il
sistema oltre le condizioni già esercitate in questa tesi.
